% Part: sets-functions-relations
% Chapter: functions
% Section: functions-relations

\documentclass[../../../include/open-logic-section]{subfiles}

\begin{document}

\olfileid{sfr}{fun}{rel}

\olsection{函数作为关系}

\begin{explain}
一个将~$A$ 中元素映射到~$B$ 中元素的函数，显然定义了~$A$ 与~$B$ 之间的一个关系，
即当且仅当 $f(x) = y$ 时 $x$ 与 $y$ 之间所具有的关系。
事实上，我们甚至——例如，若我们旨在归约数学的基本构件——可以将函数~$f$ 与这个关系\emph{等同}起来，
即等同于一个有序对的集合。
这样一来就引出了一个问题：哪些关系能以这种方式定义函数？
\end{explain}

\begin{defn}[函数的图]
设 $f\colon A \to B$ 是一个函数。$f$ 的\emph{图}是关系 $R_f \subseteq A \times B$，定义为
\[
R_f = \Setabs{\tuple{x,y}}{f(x) = y}.
\]
\end{defn}

\begin{explain}
根据外延性原理，函数的图是唯一确定的。
此外，（集合的）外延性原理也直接证实了函数隐含的外延性原理：
若 $f$ 与 $g$ 有相同的定义域和陪域，且在所有自变元上的取值都相同，则它们等同。

类似地，如果一个关系是“函数性的”，那么它就是某个函数的图。
\end{explain}

\begin{prop}\ollabel{prop:graph-function}
设 $R \subseteq A \times B$ 满足：
\begin{enumerate}
\item 若 $Rxy$ 且 $Rxz$，则 $y = z$；并且
\item 对每个 $x \in A$，都存在某个 $y \in B$ 使得 $\tuple{x, y} \in R$。
\end{enumerate}
则 $R$ 是函数 $f\colon A \to B$ 的图，其中 $f$ 由 $f(x) = y$ 当且仅当 $Rxy$ 定义。
\end{prop}

\begin{proof}
假设存在 $y$ 使得 $Rxy$。若还存在另一个 $z \neq y$ 使得 $Rxz$，则会违背关于 $R$ 的条件。
因此，如果存在 $y$ 使得 $Rxy$，那么这个 $y$ 是唯一的，所以 $f$ 是良定义的。
显然，$R_f = R$。
\end{proof}

\begin{explain}
每个函数 $f\colon A \to B$ 都有一个图，即 $A \times B$ 上由 $f(x) = y$ 定义的关系。
另一方面，每个满足 \olref{prop:graph-function} 中所述性质的关系 $R \subseteq A \times B$，都是某个函数 $f \colon A \to B$ 的图。
由于函数与其图之间的这种紧密联系，我们可以简单地将函数视作其图。
换言之，函数可以与某些关系等同，即与某些有序组的集合等同。
\oliflabeldef{sfr:rel:ref:sec}{不过请注意，这种“等同”的精神与 \olref[sfr][rel][ref]{sec} 中一致：它并非关于函数形而上学层面的断言，而是指出将函数\emph{处理}为某些集合是很方便的。这种便利性的一个原因是，现在我们可以考虑对函数执行类似于对关系所做的操作（参见 \olref[sfr][rel][ops]{sec}）。具体来说：}{}
\end{explain}

\begin{defn}\ollabel{defn:funimage}
设 $f \colon A \to B$ 是一个函数，且 $C\subseteq A$。

$f$ 在 $C$ 上的\emph{限制}是函数~$\funrestrictionto{f}{C}\colon C \to B$，其对所有 $x \in C$ 由 $(\funrestrictionto{f}{C})(x) = f(x)$ 定义。
换句话说，$\funrestrictionto{f}{C} = \Setabs{\tuple{x, y} \in R_f}{x \in
C}$。

将 $f$ \emph{应用于} $C$ 的结果是 $\funimage{f}{C} =
\Setabs{f(x)}{x \in C}$。我们也称此为 $C$ 在 $f$ 下的\emph{像}。
\end{defn}

\begin{explain}
由这些定义可知，对任意函数 $f$，$\ran{f} =
\funimage{f}{\dom{f}}$。
\oliflabeldef{sfr:rel:ops:sec}{鉴于 \olref[sfr][rel][ops]{sec} 中的定义以及我们将函数等同于关系的做法，这些概念完全符合预期。
但另外两个操作——逆和相对积——需要稍加详述。我们将在 \olref[inv]{sec} 和 \olref[cmp]{sec} 中给出详细说明。}{}
\end{explain}

\end{document}